% SOURCE_AUTHORITY: sga5_fr_workpass.tex, lines 13695--13729 (LF-normalized unit).
% SOURCE_SHA256: 791F4EFFC5E02832D5D77ED03518C8156D6F07E4C8238B03545DB93D883FBB28
\medskip

\begin{prooftext}[Demonstração]
Introduziremos o endomorfismo
\[
\psi : G \longrightarrow G
\]
definido por:
\[
\psi = \mathrm{int}_g \circ \varphi \quad , \quad \mathrm{int}_g(\xi) = g^{-1}\xi g \quad .
\]

É fácil ver que:
\begin{enumerate}
\item[i)] $\psi(G_{x'}) \subset G_{g^{-1}f'x'}$, onde $G_{x'}$ é o estabilizador do ponto $x'$.
\item[ii)] $c^\varphi(g) = \card(G^\psi)$, $G^\psi = \{ \xi \mid \psi(\xi) = \xi \}$.
\item[iii)] $G^\psi$ opera sobre ${X'_x}^{g^{-1}f'}$.
\item[iv)] $\nu_{x'}(g^{-1}f') = \nu_{hx'}(g^{-1}f')$ para cada $h \in G^\psi$.
\item[v)] $(g^{-1}f')(g_1x) = \psi(g_1)(g^{-1}f')(x)$.
\end{enumerate}

Consequentemente, pode-se decompor ${X'_x}^{g^{-1}f'}$ como a união de suas órbitas, e obtém-se:

\begin{equation}\tag{4.9.1}
\sum_{x' \in {X'_x}^{g^{-1}f'}} (\nu_{x'}(g^{-1}f') - 1) = \sum_{x' \in {X'_x}^{g^{-1}f'}/G^\psi} \frac{(\nu_{x'}(g^{-1}f') - 1)}{\card(G^\psi_{x'})} \card(G^\psi) \quad ,
\end{equation}
\[
(\nu_{x'}(g^{-1}f')=\nu\bigl((g^{-1}f')_{\mathcal{O}_{x'}}\bigr)) \quad ,
\]
Ou seja:
\begin{equation}\tag{4.9.2}
S_x(f',f,G,\varphi)(g)=\sum_{x' \in {X'_x}^{g^{-1}f'}/G^\psi}
\frac{\nu_{x'}(g^{-1}f')-1}{\card(G^\psi_{x'})} \quad ,
\end{equation}
e a proposição resulta da proposição 4.6.
\end{prooftext}
